\documentclass[10pt,a4paper]{amsart}
\usepackage[margin=19mm]{geometry}
\usepackage{amsmath,amssymb,mathtools}
\usepackage[scr=rsfs,cal=euler]{mathalfa}
\usepackage{xfrac}
\usepackage[unicode,hidelinks,bookmarks=false]{hyperref}
\newcommand{\ie}{i.e.\ }
\newcommand{\cf}{cf.\ }
\newcommand{\resp}{respectively\ }
\newcommand{\Subsection}{\subsection}
\providecommand{\nobreakdash}{\nobreak\hbox{-}}
\setlength{\emergencystretch}{2em}
\multlinegap0pt
\usepackage{mathtools}
\usepackage{amssymb}
\usepackage{float}
\usepackage{enumerate}
\usepackage{algorithm}\usepackage{algpseudocode}
\usepackage{tikz}
\newtheorem{theorem}{Theorem}
\DeclareMathOperator*{\argmin}{arg\,min}
\newcommand{\p}{\mathbb{P}}
\title{Copula Transformations for Data-Consistent Inversion}
\author{Troy Butler, Tianyi Jiang, Jo\~ao Silva, Harri Hakula, Timothy Wildey}
\date{Source: arXiv:2609.02832v2 (CC BY 4.0)}
\begin{document}
\maketitle

\noindent\emph{Source and licence.} Copula Transformations for Data-Consistent Inversion. Version arXiv:2609.02832v2 (CC BY 4.0), distributed by arXiv under the Creative Commons Attribution 4.0 International licence, http://creativecommons.org/licenses/by/4.0/, as recorded in the arXiv licence metadata for this exact version and shown on the abstract page https://arxiv.org/abs/2609.02832v2. Full licence notice: \texttt{licenses/arxiv-2609.02832-CC-BY-4.0.txt}.\par\smallskip

\noindent\textit{Editorial scope.} Counted unit: the article's theorem thm:mtd2\_converges (source0.tex lines 997--1007). The article's complete original proof of it is delivered with it (source0.tex lines 1009--1054, one \texttt{proof} environment of 46 lines). No mathematics is written, repaired or re-derived: the statement and the proof are the article's own lines, in the article's own notation, and no retained line is substituted or re-flowed.  No further source-local context is needed: every cross-reference of every retained line resolves inside the two delivered spans.  Nothing was omitted from any retained span, so the package carries no omission record.  No retained line contains a citation, so the wrapper carries no bibliography and no bibliography entry.  No retained line contains a figure environment, an image-inclusion command or drawing code, so no image is delivered and the package declares no asset dependency.  Editorial formatting only, in addition: an AMS \texttt{amsart} preamble that restates the article's own declarations and macros used by the retained lines, namely the environments \texttt{theorem} on separate counters, together with the standard packages those lines need.  The retained environments print in this document as Theorem 1 in order of appearance, whereas the article prints the same items with the numbering of its own full document.  The complete native source of the article as distributed in the arXiv e-print for this exact version is bundled unchanged as \texttt{sources/209239/source0.tex}.
\begin{theorem}\label{thm:mtd2_converges}
Let $\p_{1} \supset \p_{2} \supset \dots$, where each $\p_k$ is a nonempty convex set that is closed in the total variation topology and for which the corresponding I-projection of $P^0$ exists with finite KL divergence.
Let $\p= \bigcap_{j=1}^{\infty}\p_j$ and suppose that the I-projection of $P^0$ onto $\p$ exists with finite KL divergence and that
the sequence $\{\argmin_{P \in \p_{k}}D_{\rm KL}(P \| P^{0})\}$ is relatively compact.
Then,
$
\argmin_{P \in \p_{k}}D_{\rm KL}(P \| P^{0})
\xrightarrow{\mathrm{TV}}
\argmin_{P \in \p}D_{\rm KL}(P \| P^{0})
$.
\end{theorem}

\begin{proof}
By construction, it is clear that
\begin{equation}\label{eq:limsup}
    \limsup_{k} \inf_{P \in \p^{k}} D_{\rm KL}(P || P^{0}) \le \inf_{P \in \p} D_{\rm KL}(P || P^{0}).
\end{equation}
Since the sequence
$\left\{
\arg\min_{P\in\p_k}
D_{\mathrm{KL}}(P\|P^0)
\right\}_{k\in\mathbb{N}}$
is relatively compact in total variation, there exists a subsequence $\{k_\ell\}$ and a probability measure $P^\ast$ such that
\[
\arg\min_{P\in\p_{k_\ell}}
D_{\mathrm{KL}}(P\|P^0)
\xrightarrow{\mathrm{TV}}
P^\ast.
\]
For any fixed $j$, we have $k_\ell\geq j$ for all sufficiently large
$\ell$, and therefore, by nestedness,
\[
P_{k_\ell}^\ast\in\p_{k_\ell}\subseteq\p_j.
\]
Since $\p_j$ is closed in the total variation topology and
$P_{k_\ell}^\ast\xrightarrow{\mathrm{TV}}P^\ast$, it follows that
$P^\ast\in\p_j$. Since this holds for every $j$,
\[
P^\ast\in\bigcap_{j=1}^{\infty}\p_j=\p.
\]
By lower semicontinuity of the KL divergence with respect to total variation,
\begin{equation}\label{eq:liminf}
    \liminf_{k_{l}} \inf_{P \in \p^{k_{l}}} D_{\rm KL}(P || P^{0}) = \liminf_{k_{l}}D_{\rm KL}\left(\argmin_{P \in \p^{k_{l}}}D_{\rm KL}(P | P^{0})\, \| \, P^{0}\right)  \ge  D_{\rm KL}(P^{*} || P^{0}) \ge \inf_{P \in \p} D_{\rm KL}(P || P^{0}).
\end{equation}
By~\eqref{eq:limsup} and~\eqref{eq:liminf}, we have
\begin{equation*}
    \lim_{k_{l}} \inf_{P \in \p^{k_{l}}} D_{\rm KL}(P || P^{0}) = \inf_{P \in \p} D_{\rm KL}(P || P^{0}).
\end{equation*}
By the existence and uniqueness of the minimizer of the KL divergence and relative compactness, we have
\begin{equation*}
    \lim_{k_{l}} \argmin_{P \in \p^{k_{l}}}D_{\rm KL}(P \| P^{0}) = \argmin_{P \in \p}D_{\rm KL}(P \| P^{0}).
\end{equation*}
Since the above argument can be applied to any subsequence of $\argmin_{P \in \p^{k}}D_{\rm KL}(P | P^{0})$ to extract a further subsequence with the same limit, it follows that
\begin{equation*}
    \lim_{k} \argmin_{P \in \p^{k}}D_{\rm KL}(P \| P^{0}) = \argmin_{P \in \p}D_{\rm KL}(P \| P^{0}),
\end{equation*}
which finishes the proof.
\end{proof}

\end{document}
